\chapter{The problem, and why it is worth measuring}
\label{ch:problem}

\section{What breaks}

A browser fills a form field automatically only when the markup gives it enough
to recognise what the field is for. When it cannot, the person on the other side
of the screen types their street address by hand, on a phone, at the last step
of a checkout. That is not a cosmetic defect. It is the point in a funnel where
people stop.

The web platform already says how to avoid it. The WHATWG HTML Living Standard
defines the \code{autocomplete} attribute and a fixed set of autofill field name
tokens, and the browser vendors' own form guidance explains what their
heuristics do with them \cite{whatwg-autofill,webdev-forms,chrome-autofill}. The
hard part is not knowing the rule. The hard part is standing in front of a page
somebody else wrote, with sixty controls on it, and working out which ones are
broken and exactly what to add to each.

That is the gap this tool sits in. It loads a page in a real browser, walks the
document including same origin frames and open shadow roots, collects the
signals a browser's own heuristics would use, classifies every control against
the specification's token set, and reports the controls that will not autofill
together with the attribute to add to each one.

\section{The one design commitment everything else follows from}

The label space \emph{is} the specification's token set. Not a bespoke ontology
that is later mapped onto it.

This is worth stating early because it removes an entire class of bug from the
product. A control classified as \code{postal-code} produces the advice
\emph{add} \code{autocomplete="postal-code"}, and there is no second,
undertested translation layer between the prediction and the fix. A project that
invented its own categories would need a mapping table from category to token,
that table would be wrong somewhere, and the wrongness would surface as
confident bad advice on somebody's production markup.

\section{Why the output needs evidence rather than confidence}

The output of this tool is a list of accusations. \emph{This field will not
autofill, and here is why.} A developer acts on those, edits production markup,
and ships.

A wrong finding costs them time and costs the tool its credibility. An
unattributed finding costs more than either, because a finding that says
\emph{this looks like a postal code} without saying which signal fired and how
sure the classifier is cannot be argued with, which means it cannot be checked,
which means the only available responses are to trust it or to delete the check.

So every finding this tool emits carries a \code{signals} list naming the
concrete inputs that triggered it, and, from the point at which a probabilistic
classifier exists, a calibrated confidence. Below the reported threshold the
tool says \code{LOW\_CONFIDENCE} or \code{UNDETECTABLE\_FIELD} rather than
guessing quietly. There is no third branch, and in particular there is no path
by which a rule table's default case becomes a confident answer.

\section{What this report is}

The tool is one half of the project. The other half is the question of whether
any of it works, and that question has a specific form here.

\begin{quote}
Does a fifty kilobyte linear model that runs in microseconds on a CPU get close
enough to a twelve billion parameter language model to be the right default for
a developer tool?
\end{quote}

That is the headline experiment. It is a product question rather than a research
one, and it was chosen because it is the question an engineer actually faces
when deciding what to put in a wheel that people install with one command. This
report answers it with numbers, and prints the answer the measurement gave
rather than the one the project was built expecting.

The measurement itself is described in Chapter~\ref{ch:method}, the numbers are
in Chapter~\ref{ch:headline}, the significance procedure and its verdicts are in
Chapter~\ref{ch:statistics}, and the predictions that were registered before any
of it ran are scored in Chapter~\ref{ch:predictions}. What the answer does not
license is in Chapter~\ref{ch:limitations}, and it is a longer chapter than the
answer.

\section{How to read the numbers in this document}

\syntheticnote

Two properties of this document are worth knowing before the first table.

\textbf{No number here was typed by hand.} Every table is generated from a
committed result file by \code{scripts/make\_report\_tables.py}, every figure by
\code{scripts/make\_report\_figures.py}, and every number quoted in a sentence
is a macro defined by the first of those beside the path it was read from. The
chapters of this report contain no digits at all. A continuous integration job
scans these sources for anything shaped like a measurement and fails the build
unless the line names the result file behind it, then follows that file to its
run manifest and the manifest to a commit object this repository can produce.

\textbf{Every result cited here was produced from a clean working tree.} A run
taken from a tree with uncommitted changes is marked dirty in its manifest and
is not citable, and the traceability job fails on a citation of one. One
artefact in this repository does carry that mark, and Chapter~\ref{ch:method}
says which and what follows from it, rather than leaving the reader to find it.
